
\documentclass{article}
\usepackage{colortbl}
\usepackage{makecell}
\usepackage{multirow}
\usepackage{supertabular}

\begin{document}

\newcounter{utterance}

\centering \large Interaction Transcript for game `bbh', experiment `object\_counting', episode 59 with reflexion{-}Qwen3.5{-}27B.
\vspace{24pt}

{ \footnotesize  \setcounter{utterance}{1}
\setlength{\tabcolsep}{0pt}
\begin{supertabular}{c@{$\;$}|p{.15\linewidth}@{}p{.15\linewidth}p{.15\linewidth}p{.15\linewidth}p{.15\linewidth}p{.15\linewidth}}
    \# & \multicolumn{2}{c}{Player} && \multicolumn{2}{c}{Game Master} \\
    \hline

    \theutterance \stepcounter{utterance}  
    & & & \multicolumn{4}{p{0.6\linewidth}}{
        \cellcolor[rgb]{0.9,0.9,0.9}{
            \makecell[{{p{\linewidth}}}]{
                \texttt{\tiny{[P1$\langle$GM]}}
                \texttt{You have attempted this same task before and did not succeed. These are the plans you wrote after those attempts.} \\
\texttt{Trial \#1: New plan: My previous attempt failed because I included a verbose explanation and a step{-}by{-}step breakdown before providing the final answer, whereas the few{-}shot examples demonstrated that the expected response is strictly the final number. The pattern in the examples (Q: ... A: {[}Number{]}) indicates that no intermediate reasoning or text should be generated. Next time, I will immediately output only the calculated integer (in this case, 8) without any preceding text, lists, or the "A:" prefix, to strictly match the few{-}shot format.} \\
\\ 
\\ 
\texttt{Questions that involve enumerating objects and asking the model to count them.} \\
\\ 
\texttt{Q: I have a blackberry, a clarinet, a nectarine, a plum, a strawberry, a banana, a flute, an orange, and a violin. How many fruits do I have?} \\
\texttt{A: 6} \\
\\ 
\texttt{Q: I have an orange, a raspberry, two peaches, a blackberry, an apple, a grape, a nectarine, and three plums. How many fruits do I have?} \\
\texttt{A: 11} \\
\\ 
\texttt{Q: I have a lettuce head, a head of broccoli, an onion, a stalk of celery, two carrots, a garlic, and a yam. How many vegetables do I have?} \\
\texttt{A: 8} \\
\\ 
\texttt{Q: I have two beds, a stove, a microwave, a table, a fridge, a lamp, and a toaster. How many objects do I have?} \\
\texttt{A:} \\
            }
        }
    }
     \\ \\

    \theutterance \stepcounter{utterance}  
    & \multicolumn{4}{p{0.6\linewidth}}{
        \cellcolor[rgb]{0.561,0.737,0.561}{
            \makecell[{{p{\linewidth}}}]{
                \texttt{\tiny{[P1$\rangle$GM]}}
                \texttt{8} \\
            }
        }
    }
    & & \\ \\

    \theutterance \stepcounter{utterance}  
    & & & \multicolumn{2}{p{0.3\linewidth}}{
        \cellcolor[rgb]{0.9,0.9,0.9}{
            \makecell[{{p{\linewidth}}}]{
                \texttt{\tiny{[GM$|$GM]}}
                \texttt{Target: 8} \\
            }
        }
    }
    & & \\ \\

    \theutterance \stepcounter{utterance}  
    & & & \multicolumn{2}{p{0.3\linewidth}}{
        \cellcolor[rgb]{0.9,0.9,0.9}{
            \makecell[{{p{\linewidth}}}]{
                \texttt{\tiny{[GM$|$GM]}}
                \texttt{game\_result = WIN} \\
            }
        }
    }
    & & \\ \\

\end{supertabular}
}

\end{document}
